% !PRINTED 461
\weberpage{461}
%<<CONT
hence
\[
-\gamma_2\left(\frac{-3+\omega}{2}\right) = f(\omega)^8 - \frac{256}{f(\omega)^{16}}.
\]

If one therefore sets:
\weq{20}{q = e^{-\pi\sqrt{m}},}
one obtains [\S\,24, (11)]:
\[
Z = -\gamma_2\left(\frac{-3+i\sqrt{m}}{2}\right)
\]
\[
= q^{-\frac{1}{3}} \overset{\nu}{\underset{1,\infty}{\Pi}} (1+q^{2\nu-1})^8 - \frac{256\,q^{\frac{2}{3}}}{\overset{\nu}{\underset{1,\infty}{\Pi}} (1+q^{2\nu-1})^{16}}.
\]

We now make use of the determination of limits valid for every properly fractional positive $x$:
\[
1+x < e^x < \frac{1}{1-x}, \quad e^{-x} > 1-x,
\]
and obtain:
\[
1 < \Pi(1+q^{2\nu-1}) < e^{\frac{q}{1-q^2}},
\]
from which
\[
q^{-\frac{1}{3}} - 256\,q^{\frac{2}{3}} < Z < q^{-\frac{1}{3}} e^{\frac{8q}{1-q^2}} - 256\,q^{\frac{2}{3}} e^{-\frac{16q}{1-q^2}},
\]
and by enlarging the upper limit still further, one may also put in its place:
\[
Z < q^{-\frac{1}{3}} - 256\,q^{\frac{2}{3}} + \frac{8\,q^{\frac{2}{3}}}{1-8\,q-q^2} + \frac{2^{12}\,q^{\frac{5}{3}}}{1-q^2}.
\]
For $m=11$, $q$ is approximately $=2^{-15}$, from which one sees that the difference of the two limits:
\[
\frac{8\,q^{\frac{2}{3}}}{1-8\,q-q^2} + \frac{2^{12}\,q^{\frac{5}{3}}}{1-q^2}
\]
is already very small for $m=11$, and therefore still more so for the larger values of $m$. $Z$ is therefore the integer lying next above
\[
q^{-\frac{1}{3}} - 256\,q^{\frac{2}{3}}
\]
and this value, for the larger $m$ already $q^{-\frac{1}{3}}$, come extraordinarily close to an integer. % [sic] German print has singular "dieser Wert" with plural "kommen"
From this one computes these numbers\footnote{One does well to fix in mind, once and for all, for these and for many of the calculations described later, the \gesperrt{Briggian} logarithm
\[
\log(\pi \log e) = 0{,}134\,934\,184\,0
\]
for which, moreover, 7 decimal places are usually sufficient.}:
%CONT>>
